\section{Nonvanishing and the rank-weighted tail}\label{sec:tail}

We now combine the transferred model estimates, the integer product
event, and the terminal mean square to count large orders of vanishing.
A separate bound for each individual rank then converts this count into
the rank-weighted tail. The companion density theorem enters only in
the final subsection, after the tail theorem is proved.
The constants and the fixed
lower threshold \(k_0\) are increased finitely many times in this section.
All choices depend only on the fixed family and the already fixed
smoothness orders.

\Needspace{11\baselineskip}
\subsection{Counting large orders of vanishing}

\begin{proposition}\label{tail:count}
There are \(k_0\), \(X_0\), and \(C\), depending on \(E\), such that for
every \(F\in\mathcal F\), every dyadic \(X\ge X_0\), and every integer
\[
 k_0+1\le k\le(\log X)^{3/5},
\]
one has
\[
 \#\{d\in\mathcal A(X):a(F^{(d)})>k\}\le C Xk^{-2}.
\]
\end{proposition}

\begin{proof}
Write \(\ell=\log X\), and first consider the functional-equation sign
\((-1)^k\), with \(k_0\le k\le\ell^{3/5}\).
Fix \(c=\pm1\), and use the weights and partition of
Section~\ref{sec:weights}. We write \(P_j^{\rm eu}\) for the positive
Euler products of Section~\ref{sec:model}.

Corollary~\ref{weights:comparison-interface} and
Lemma~\ref{model:mollifier-length} permit
Proposition~\ref{analytic:comparison} to be applied to \(H_jS_j\),
for \(0\le j\le J\), with \(\Delta=\delta_j/4\).
By the first assertion of Proposition~\ref{model:main-model-estimates},
outside \(O(X/k^2)\) integers in \(X\le|m|<2X\),
\begin{equation}\label{eq:tail-main-good}
 |H_0(z_m)S_{0,z_m}(m/X)-A_{m/X}|<\frac12 A_{m/X}.
\end{equation}
Indeed \(A_x=(C_F|x|)^s\Gamma(1+s)\) is bounded above and below by positive
constants on the fixed \(x\)-range, uniformly for \(s=k/\ell\le\ell^{-2/5}\).
Chebyshev's inequality applies to the transferred \(O(k^{-2})\)
mean-square error. The exponentially small comparison error is absorbed
uniformly.

For \(1\le j\le J\), the second model estimate gives
\[
 \frac1X\sum_m\Psi(m/X)|H_j(z_m)S_{j,z_m}(m/X)|^2
 \ll (k2^j)^{2C_5}2^{-2jk}+e^{-c\delta_j\ell}
 \ll 2^{-3jk/2}.
\]
The last inequality follows first by increasing fixed \(k_0\) to absorb
the polynomial factor. For the comparison error, use
\[
 \delta_j\ell\ge\ell^{9/10},\qquad
 jk\ll\ell^{3/5}\log\ell,
\]
and then increase one \(X_0\) for all the permitted \(k,j\).
The terminal estimate of Lemma~\ref{weights:terminal-small} gives
\[
 \frac1X\sum_{d\in\mathcal A(X)}
       |S_{*,\chi_d}(d/X)|^2\ll 2^{-3Jk/2}.
\]
Chebyshev and summation of the resulting geometric probabilities imply,
outside \(O(X/k^2)\) further admissible parameters,
\begin{equation}\label{eq:tail-other-good}
 |H_j(\chi_d)S_{j,\chi_d}(d/X)|\le2^{-jk/2}\quad(1\le j\le J),
 \qquad
 |S_{*,\chi_d}(d/X)|\le2^{-Jk/2}.
\end{equation}
Here \(\sum_{j\ge1}2^{-jk/2}\ll k^{-2}\) after fixing \(k_0\), and \(J\ge1\)
for sufficiently large \(X\).

Intersect these events with the event in
Proposition~\ref{model:good-event}. Its additional exceptional set is
also \(O(X/k^2)\). On the intersection,
\[
 \frac12\le H_jP_j^{\rm eu}\le2,\qquad
 P_j^{\rm eu}\le 2^{jk/8}P_0^{\rm eu},\qquad
 1\le 2^{Jk/8}P_0^{\rm eu}.
\]
Choose a fixed \(a_*>0\) with \(A_x\ge a_*\) for all the parameters.
Equation~\eqref{eq:tail-main-good} yields
\[
 |S_{0,\chi_d}(d/X)|\ge\frac{a_*}{4}P_0^{\rm eu}.
\]
Equation~\eqref{eq:tail-other-good} yields
\[
 \sum_{j=1}^J|S_{j,\chi_d}(d/X)|+|S_{*,\chi_d}(d/X)|
 \le P_0^{\rm eu}
     \left(2\sum_{j=1}^J2^{-3jk/8}+2^{-3Jk/8}\right).
\]
For a sufficiently large fixed \(k_0\), the last parenthesis is less
than \(a_*/4\). The entire weighted sum is therefore nonzero.
Lemma~\ref{weights:vanishing} excludes the simultaneous conditions
\(a(F^{(d)})>k\) and \(\epsilon_{F^{(d)}}=(-1)^k\).
Including both values of \(c\), we have proved
\begin{equation}\label{eq:tail-parity-count}
 \#\{d\in\mathcal A(X):
       a(F^{(d)})>k,\ \epsilon_{F^{(d)}}=(-1)^k\}
       \ll Xk^{-2}.
\end{equation}
For the other functional-equation sign, apply
\eqref{eq:tail-parity-count} at \(k-1\). A rank exceeding \(k\) also exceeds
\(k-1\), and \((k-1)^{-2}\ll k^{-2}\).
This proves the proposition.
\end{proof}

\subsection{A uniform bound for an individual rank}

\begin{lemma}\label{tail:jensen}
Uniformly for \(F\in\mathcal F\) and \(d\in\mathcal A(X)\),
\[
 a(F^{(d)})\ll_E\log(3X).
\]
\end{lemma}

\begin{proof}
Put \(L=L_{F^{(d)}}\) and \(C_d=C_F|d|\).
On \(\Re s=5\), absolute convergence and
\eqref{eq:setup-coefficients} give a uniform bound.
On \(\Re s=-1\), the functional equation is
\[
 L(s)=\epsilon_{F^{(d)}} C_d^{\,1-2s}
          \frac{\Gamma(3/2-s)}{\Gamma(1/2+s)}L(1-s).
\]
At \(s=-1+it\), the modulus of the gamma quotient equals
\[
 |(-1/2+it)(1/2+it)(3/2+it)|.
\]
Since \(L(2-it)\) is uniformly bounded,
\[
 |L(-1+it)|\ll_E X^3(1+|t|)^3.
\]
Apply the strip principle to \(L(s)/(s+2)^4\) on
\(-1\le\Re s\le5\). To see that the possibly conductor-dependent
finite-order growth in the strip introduces no new constant, first
multiply by
\[
 \exp\!\left(-\varepsilon
       \cos\frac{\pi(s-2)}{12}\right),\qquad \varepsilon>0.
\]
Its modulus is at most one on this strip and decays faster than any
finite-order growth as \(|\Im s|\to\infty\).
The maximum principle on expanding rectangles gives the bound from
the two vertical sides. Letting \(\varepsilon\) decrease to zero shows
that \(L(s)/(s+2)^4\ll_E X^3\) throughout the strip.
In particular, the maximum of \(|L|\) on the closed disk
\(|s-2|\le5/2\) is \(O_E(X^3)\).

At the center, the Euler factors and their parameter bounds give
\[
 |L(2)|\ge \prod_p(1+p^{-2})^{-2}
          =\left(\frac{\zeta(4)}{\zeta(2)}\right)^2>0.
\]
Jensen's formula, with inner radius \(2\) and outer radius \(5/2\),
therefore bounds the number of zeros in the inner disk, counted with
multiplicity, by \(O_E(\log(3X))\).
The point \(1/2\) lies in that inner disk. This proves the result.
\end{proof}

\subsection{Summing the integer tails}

\begin{proof}[Proof of Theorem~\ref{thm:tail}]
We first work in \(\mathcal A(X)\) for a fixed \(F\in\mathcal F\).
Let
\[
 T_{F,X}(k)=\#\{d\in\mathcal A(X):a(F^{(d)})>k\},
 \qquad K=\lfloor(\log X)^{3/5}\rfloor.
\]
For an integer \(R\ge k_0+1\) and \(X\) large enough that \(K>R\),
the elementary integer-tail bound gives
\begin{align}
 \sum_{\substack{d\in\mathcal A(X)\\a(F^{(d)})>R}}a(F^{(d)})
 &\le R T_{F,X}(R)
       +\sum_{k=R}^{K-1}T_{F,X}(k)
       +C_E\log X\,T_{F,X}(K) \notag\\
 &\ll_E \frac X R+X(\log X)^{-1/5}.
 \label{eq:tail-dyadic-mass}
\end{align}
For ranks between \(R+1\) and \(K\), the first two terms count the rank
exactly. For ranks exceeding \(K\), the last term bounds the remaining
mass by Lemma~\ref{tail:jensen}.
Proposition~\ref{tail:count} supplies every counting estimate used here.
In particular the last term is
\(O_E(X\log X/K^2)=O_E(X(\log X)^{-1/5})\).
The constant in \eqref{eq:tail-dyadic-mass} is independent of \(R\).

For completeness, fix \(R\) and cover \(0<|d|\le Y\) by dyadic blocks.
Enlarge the last partial block to a full block, using nonnegativity.
The sum of their left endpoints is less than \(2Y\).
The finitely many blocks below the threshold for
\eqref{eq:tail-dyadic-mass} contribute a finite number \(B_{E,R}\).
For the error term in the remaining blocks, split at \(X=\sqrt Y\).
The smaller blocks contribute \(O_E(\sqrt Y)\), and in the larger ones
\(\log X\ge\frac12\log Y\).
Consequently
\[
 \sum_{\substack{d\ {\rm admissible}\\0<|d|\le Y\\a(F^{(d)})>R}}
     a(F^{(d)})
 \le B_{E,R}+\frac{C_EY}{R}
       +O_E\bigl(\sqrt Y+Y(\log Y)^{-1/5}\bigr).
\]
After division by \(Y\) and passage to the limsup, the desired
\(C_E/R\) bound holds for each \(F\in\mathcal F\).
Lemma~\ref{setup:reduction}, its bounded multiplicity, and the finite
maximum of these constants prove the same bound for every original
signed squarefree parameter of \(E\).
\end{proof}

\subsection{The mean}

\begin{proof}[Proof of Theorem~\ref{thm:mean}]
Write \(N(Y)=\#\mathcal D(Y)\) and let \(N_j(Y)\) count the parameters
of analytic rank \(j\).
The elementary squarefree sieve gives
\[
 N(Y)=2\sum_{n\le Y}\mu(n)^2
     =\frac{2Y}{\zeta(2)}+O(\sqrt Y).
\]
Indeed, expand \(\mu(n)^2=\sum_{r^2\mid n}\mu(r)\), interchange the
finite sums, and bound the omitted tail of
\(\sum_r\mu(r)/r^2=1/\zeta(2)\).

By the analytic density theorem \eqref{eq:density-input},
\(N_0(Y)/N(Y)\to1/2\) and \(N_1(Y)/N(Y)\to1/2\).
Thus \(N(Y)-N_0(Y)-N_1(Y)=o(Y)\).
For every fixed integer \(R\ge R_E\),
\[
 \sum_{\substack{d\in\mathcal D(Y)\\2\le a(E^{(d)})\le R}}
        a(E^{(d)})
 \le R\bigl(N(Y)-N_0(Y)-N_1(Y)\bigr)=o_{E,R}(Y).
\]
Theorem~\ref{thm:tail} now implies
\[
 \limsup_{Y\to\infty}\frac1Y
 \sum_{\substack{d\in\mathcal D(Y)\\a(E^{(d)})\ge2}}a(E^{(d)})
 \le\frac{C_E}{R}.
\]
Let \(R\to\infty\). The left side is zero by nonnegativity.
The total rank sum is consequently \(N_1(Y)+o_E(Y)\).
Dividing by \(N(Y)\) proves the limit \(1/2\).
\end{proof}
